% Archive record 48, line 1.
\hypertarget{part-3-48-0001-iii-f-rsi-failure-modes-experimental-instrumentation}{%
\section{III-F --- RSI Failure Modes \& Experimental Instrumentation}\label{part-3-48-0001-iii-f-rsi-failure-modes-experimental-instrumentation}}

这是 Part III 的最后一个模块，也是目前最接近实际科研工作的一个模块。

此前我们已经理解了 RSI 的统计机制、系统架构、搜索算法以及代表性研究论文。现在假设导师或博士生直接给我们一个能够运行的 Financial RSI Framework，我们需要回答：

\textbf{如何将一条看似简单的 IS/OOS Performance Curve，转化为具有科学解释力的研究结论？}

这实际上涉及四个层次的研究工作。

\textbf{01}

\textbf{Observation --- 观察}

记录每个 Agent 版本的行为和性能。

\textbf{02}

\textbf{Detection --- 检测}

判断性能是否真的退化，以及何时开始退化。

\textbf{03}

\textbf{Diagnosis --- 诊断}

识别可能造成退化的 Revision 和统计机制。

\textbf{04}

\textbf{Intervention --- 干预}

通过受控实验验证缓解机制是否有效。

这四个层次恰好对应 X Program \#161 的核心研究任务。

本阶段会减少基础概念的重复介绍，重点学习新的实验分析方法。

\hypertarget{part-3-48-0001-rsi-failure-taxonomyux5931ux8d25ux673aux5236ux5206ux7c7b}{%
\subsection{1. RSI Failure Taxonomy（失败机制分类）}\label{part-3-48-0001-rsi-failure-taxonomyux5931ux8d25ux673aux5236ux5206ux7c7b}}

首先，我们将 Part II 和 Part III 中已经讨论过的失败机制统一整理。

\begin{readertable}
\begin{longtable}[]{@{}
  >{\raggedright\arraybackslash}p{(\columnwidth - 4\tabcolsep) * \real{0.30}}
  >{\raggedright\arraybackslash}p{(\columnwidth - 4\tabcolsep) * \real{0.43}}
  >{\raggedright\arraybackslash}p{(\columnwidth - 4\tabcolsep) * \real{0.23}}@{}}
\toprule
Failure Mode & 主要机制 & Knowledge Index \\
\midrule
\endhead
Classical Overfitting & 金融模型过度拟合训练数据 & II-A / B \\
Selection Bias & 从大量候选中选择受噪声影响的最优者 & II-C \\
Adaptive Overfitting & 反复利用历史评估反馈 & II-D \\
Data Leakage & 使用不应获得的信息 & II-E \\
Distribution Shift & 数据生成机制发生变化 & II-E \\
Memory-mediated Bias & 将错误经验持久化 & III-B \\
Loss of Improvability & 修改损害后续自我改进能力 & III-C \\
Premature Convergence & 过早集中于某条搜索路径 & III-D \\
Reward Hacking & 利用代理指标或评价机制的缺陷 & III-D \\
Implementation Failure & 代码错误、工具失败或功能退化 & III-A \\
\bottomrule
\end{longtable}
\end{readertable}

这十种 Failure Modes 可以进一步归纳为四个研究层面。

\textbf{Statistical Failure}

过拟合、选择偏差、评估噪声

\textbf{Optimization Failure}

局部最优、短视搜索、改进能力退化

\textbf{Evaluation Failure}

数据泄漏、指标失效、Reward Hacking

\textbf{System Failure}

代码错误、状态污染、资源耗尽

另外，Distribution Shift 属于外部环境层面的变化，可能与上述四类机制共同作用。

这里的分类是我们为本项目建立的工作性 Taxonomy，不是互斥的正式学术分类。

尤其需要注意：\textbf{观察到 OOS Decay 并不意味着已经观察到了 Overfitting。} 前者是实验现象，后者是需要进一步验证的机制解释。

\hypertarget{part-3-48-0001-experimental-instrumentationux5b9eux9a8cux89c2ux6d4bux4f53ux7cfb}{%
\subsection{2. Experimental Instrumentation（实验观测体系）}\label{part-3-48-0001-experimental-instrumentationux5b9eux9a8cux89c2ux6d4bux4f53ux7cfb}}

\hypertarget{part-3-48-0001-ux4ec0ux4e48ux662f-instrumentation}{%
\subsubsection{2.1 什么是 Instrumentation？}\label{part-3-48-0001-ux4ec0ux4e48ux662f-instrumentation}}

在软件工程和实验研究中，Instrumentation 指在系统内部设置观测点，使原本不可见的运行过程能够被记录、测量和分析。

对于 RSI，我们不能只保存：

\begin{verbatim}
Round 1: Score = 0.65
Round 2: Score = 0.71
Round 3: Score = 0.75
\end{verbatim}

因为这样的记录只能说明得分发生了变化，却无法解释变化的原因。

我们需要保存足够的信息，以重建整个 Agent Evolution Process。

\bookgraphic{0.432}{part-3-48-0001-003-01}{Instrumentation 将 Agent 的内部演化过程转化为可以进行统计分析的研究数据。}

\hypertarget{part-3-48-0001-four-levels-of-experiment-logging}{%
\subsubsection{2.2 Four Levels of Experiment Logging}\label{part-3-48-0001-four-levels-of-experiment-logging}}

建议将实验数据分成四种粒度，而不是全部混合在一个 CSV 文件中。

\begin{readertable}
\begin{longtable}[]{@{}
  >{\raggedright\arraybackslash}p{(\columnwidth - 4\tabcolsep) * \real{0.22}}
  >{\raggedright\arraybackslash}p{(\columnwidth - 4\tabcolsep) * \real{0.32}}
  >{\raggedright\arraybackslash}p{(\columnwidth - 4\tabcolsep) * \real{0.42}}@{}}
\toprule
Logging Level & 记录对象 & Key Variables \\
\midrule
\endhead
Run-level & 一次完整 RSI 实验 & Seed、LLM、Dataset、Budget \\
Iteration-level & 一轮改进过程 & Incumbent、IS/OOS Score、Query Count \\
Candidate-level & 每一个候选版本 & Parent、Code Diff、Acceptance、Failure \\
Task-level & 每个版本在具体任务上的表现 & Loss、Return、Prediction、Runtime \\
\bottomrule
\end{longtable}
\end{readertable}

这里最容易被忽略的是 \textbf{Candidate-level Logging}。

假设 Agent 在 100 次尝试中只接受了 7 次修改。如果我们只保存这 7 个成功版本，就无法充分分析 Selection Bias，因为另外 93 个被拒绝的候选同样参与了搜索与选择过程。

这实际上是一种实验记录层面的 Survivorship Bias。

\hypertarget{part-3-48-0001-recommended-data-schema}{%
\subsubsection{2.3 Recommended Data Schema}\label{part-3-48-0001-recommended-data-schema}}

一个初步的研究数据结构可以如下设计：

\begin{readertable}
\begin{longtable}[]{@{}
  >{\raggedright\arraybackslash}p{(\columnwidth - 2\tabcolsep) * \real{0.40}}
  >{\raggedright\arraybackslash}p{(\columnwidth - 2\tabcolsep) * \real{0.56}}@{}}
\toprule
File & Purpose \\
\midrule
\endhead
\texttt{runs.jsonl} & 保存每次独立实验的配置 \\
\texttt{iterations.parquet} & 保存每轮演化轨迹 \\
\texttt{candidates.parquet} & 保存全部候选及选择结果 \\
\texttt{evaluations.parquet} & 保存任务级别的评估结果 \\
\texttt{revisions.jsonl} & 保存代码修改及分类标签 \\
\texttt{versions/} & 保存完整 Agent 版本 \\
\texttt{manifest.json} & 保存数据、代码和环境的版本信息 \\
\bottomrule
\end{longtable}
\end{readertable}

例如，一条 Candidate Record 可以采用如下结构：

\begin{verbatim}
{
    "run_id": "run_001",
    "iteration": 12,
    "candidate_id": "candidate_047",
    "parent_version_id": "agent_008",
    "improver_version_id": "agent_008",
    "code_commit": "example_hash",
    "revision_categories": [
        "feature_expansion",
        "search_budget_expansion"
    ],
    "is_score": 0.812,
    "accepted": true,
    "evaluation_calls": 15,
    "input_tokens": 12500,
    "output_tokens": 3200,
    "execution_status": "success"
}
\end{verbatim}

OOS Audit Record 则由外部 Evaluator 单独维护，不能被正在优化的 Agent 读取或修改。

上述字段和目录是候选设计，并非课题组现有的数据格式。

\hypertarget{part-3-48-0001-trajectory-analysisux6f14ux5316ux8f68ux8ff9ux5206ux6790}{%
\subsection{3. Trajectory Analysis（演化轨迹分析）}\label{part-3-48-0001-trajectory-analysisux6f14ux5316ux8f68ux8ff9ux5206ux6790}}

\hypertarget{part-3-48-0001-rsi-ux5b9eux9645ux4e0aux5b58ux5728ux4e09ux79cdux4e0dux540cux7684ux65f6ux95f4ux5c3aux5ea6}{%
\subsubsection{3.1 RSI 实际上存在三种不同的时间尺度}\label{part-3-48-0001-rsi-ux5b9eux9645ux4e0aux5b58ux5728ux4e09ux79cdux4e0dux540cux7684ux65f6ux95f4ux5c3aux5ea6}}

这是一个此前没有充分展开的研究问题。

\begin{readertable}
\begin{longtable}[]{@{}
  >{\raggedright\arraybackslash}p{(\columnwidth - 4\tabcolsep) * \real{0.30}}
  >{\raggedright\arraybackslash}p{(\columnwidth - 4\tabcolsep) * \real{0.18}}
  >{\raggedright\arraybackslash}p{(\columnwidth - 4\tabcolsep) * \real{0.49}}@{}}
\toprule
Time Axis & Notation & 含义 \\
\midrule
\endhead
Proposal Index & \(j\) & 已经产生了多少候选版本 \\
Improvement Depth & \(k\) & 当前版本经历了多少次谱系上的改进 \\
Evaluation Budget & \(Q_k\) & 累计执行了多少次评估 \\
\bottomrule
\end{longtable}
\end{readertable}

假设一个系统一共生成了 100 个候选 Agent，但只有 10 个被接受。

那么：
\[\text{Proposal Count}=100\]

并不意味着：
\[\text{Improvement Depth}=100.\]

而对于 Archive-based Evolution，即使已经接受了 10 个不同版本，最终版本的谱系深度也可能小于 10。

\bookgraphic{0.477}{part-3-48-0001-004-01}{虽然 A₆ 是第六个新生成的 Agent，但它的谱系深度只有 3。}

因此，研究时应当明确区分：

\begin{itemize}
\tightlist
\item
  Chronological Trajectory：按照版本生成时间排列；
\item
  Incumbent Trajectory：按照每一时刻被选中的当前最优版本排列；
\item
  Lineage Trajectory：按照父子继承关系排列。
\end{itemize}

如果采用 Archive-based RSI，这三条曲线可能具有完全不同的含义。

\hypertarget{part-3-48-0001-core-trajectory-metrics}{%
\subsubsection{3.2 Core Trajectory Metrics}\label{part-3-48-0001-core-trajectory-metrics}}

沿用 Part I 中的记号：
\[S_k^{\mathrm{IS}},
\qquad
S_k^{\mathrm{OOS}},\]

并统一采用 Score 越高越好的定义。

\hypertarget{part-3-48-0001-metric-1-generalization-gap}{%
\paragraph{Metric 1 --- Generalization Gap}\label{part-3-48-0001-metric-1-generalization-gap}}

\[\boxed{
G_k=S_k^{\mathrm{IS}}-S_k^{\mathrm{OOS}}
}\]

衡量开发得分与样本外得分之间的差距。

\hypertarget{part-3-48-0001-metric-2-incremental-improvement}{%
\paragraph{Metric 2 --- Incremental Improvement}\label{part-3-48-0001-metric-2-incremental-improvement}}

\[\begin{aligned}
\Delta_k^{\mathrm{IS}}
&=S_k^{\mathrm{IS}}-S_{k-1}^{\mathrm{IS}},\\
\Delta_k^{\mathrm{OOS}}
&=S_k^{\mathrm{OOS}}-S_{k-1}^{\mathrm{OOS}}.
\end{aligned}\]

衡量相邻版本之间的变化。

对于树状演化结构，应当将 \(A_{k-1}\) 替换为对应的 Parent Version，才能分析实际 Revision 的影响。

\hypertarget{part-3-48-0001-metric-3-peak-round}{%
\paragraph{Metric 3 --- Peak Round}\label{part-3-48-0001-metric-3-peak-round}}

\[\boxed{
k_{\mathrm{peak}}
=
\arg\max_{0\leq k\leq K}
S_k^{\mathrm{OOS}}
}\]

记录观察到的 OOS Performance Peak。

\hypertarget{part-3-48-0001-metric-4-peak-to-final-decay}{%
\paragraph{Metric 4 --- Peak-to-Final Decay}\label{part-3-48-0001-metric-4-peak-to-final-decay}}

\[\boxed{
D_{\mathrm{PF}}
=
S_{k_{\mathrm{peak}}}^{\mathrm{OOS}}
-
S_K^{\mathrm{OOS}}
}\]

衡量最终版本相对于历史峰值的退化程度。

\hypertarget{part-3-48-0001-metric-5-endpoint-improvement}{%
\paragraph{Metric 5 --- Endpoint Improvement}\label{part-3-48-0001-metric-5-endpoint-improvement}}

\[\boxed{
I_{\mathrm{end}}
=
S_K^{\mathrm{OOS}}-S_0^{\mathrm{OOS}}
}\]

衡量 RSI 的最终净收益。

\hypertarget{part-3-48-0001-metric-6-maximum-performance-drawdown}{%
\paragraph{Metric 6 --- Maximum Performance Drawdown}\label{part-3-48-0001-metric-6-maximum-performance-drawdown}}

借鉴金融领域 Drawdown 的思想，可以定义：
\[\boxed{
D_{\max}
=
\max_{0\leq i\leq j\leq K}
\left(
S_i^{\mathrm{OOS}}-S_j^{\mathrm{OOS}}
\right)
}\]

它与 Peak-to-Final Decay 的区别在于，后者只关注最终时刻，而前者能够捕捉中途发生、后来又得到恢复的严重性能退化。

\textbf{RSI Trajectory --- Metric Illustration}

\bookgraphic{0.786}{part-3-48-0001-005-01}{}

Peak Round

\textbf{4}

Endpoint Improvement

\textbf{+0.15}

Peak-to-Final Decay

\textbf{0.15}

Maximum Drawdown

\textbf{0.28}

Final IS/OOS Gap

\textbf{0.39}

Peak OOS Performance

\textbf{0.70}

合成教学数据。最终版本虽然优于初始版本，但已经经历了严重的中途退化。

如果研究人员只报告 Endpoint Improvement，就可能将上述实验视为一次成功的 RSI。

而从完整轨迹来看，该系统显然还存在需要研究的 Generalization Problem。

需要注意，这些基于经验最大值的指标都可能受到 Extreme Value Selection Bias 的影响，因此应当结合统计不确定性进行解释。

\hypertarget{part-3-48-0001-change-point-detectionux53d8ux70b9ux68c0ux6d4b}{%
\subsection{4. Change-point Detection（变点检测）}\label{part-3-48-0001-change-point-detectionux53d8ux70b9ux68c0ux6d4b}}

这是整个 III-F 中最重要的新数学内容。

项目描述明确要求识别：

\begin{quote}
The iteration at which out-of-sample performance begins to decay.
\end{quote}

但如何定义这个时刻，并不是一件完全 Trivial 的事情。

\hypertarget{part-3-48-0001-peak-detection-ux4e0dux7b49ux4e8e-change-point-detection}{%
\subsubsection{4.1 Peak Detection 不等于 Change-point Detection}\label{part-3-48-0001-peak-detection-ux4e0dux7b49ux4e8e-change-point-detection}}

如果我们直接使用：
\[k_{\mathrm{peak}}
=
\arg\max_k S_k^{\mathrm{OOS}},\]

得到的只是观察得分最高的时刻。

由于 Evaluation Noise 的存在，这个时刻不一定是真正的性能变化点。

我们将观察得分写为：
\[S_k^{\mathrm{OOS}}
=
J_k^{\mathrm{OOS}}+\eta_k,\]

其中 \(J_k^{\mathrm{OOS}}\) 表示潜在的真实性能，\(\eta_k\) 表示评估误差。

研究者真正希望识别的是 \(J_k^{\mathrm{OOS}}\) 的趋势何时发生变化，而不是简单寻找 \(S_k^{\mathrm{OOS}}\) 的最大值。

\bookgraphic{0.477}{part-3-48-0001-006-01}{}

Observed OOS

Latent performance

合成示意：经验峰值出现在第 4 轮，而潜在趋势在第 6 轮才发生变化。

\hypertarget{part-3-48-0001-segmented-regressionux5206ux6bb5ux56deux5f52}{%
\subsubsection{4.2 Segmented Regression（分段回归）}\label{part-3-48-0001-segmented-regressionux5206ux6bb5ux56deux5f52}}

我们可以利用你已经熟悉的 Linear Regression，建立一个具有未知变点的模型。

沿用 Part I 中的参数记号：
\[\boxed{
S_k^{\mathrm{OOS}}
=
\alpha+\beta k
+\gamma(k-\tau)_+
+\eta_k
}\]

其中：
\[(k-\tau)_+=\max(0,k-\tau).\]

这里的 \(\alpha\) 为回归截距，\(\tau\) 为候选变点，并非此前 Softmax Selection 的温度参数。

展开可得：
\[S_k^{\mathrm{OOS}}
=
\begin{cases}
\alpha+\beta k+\eta_k,
&k\leq\tau,\\
\alpha+\beta k+\gamma(k-\tau)+\eta_k,
&k>\tau.
\end{cases}\]

因此：

\begin{readertable}
\begin{longtable}[]{@{}
  >{\raggedright\arraybackslash}p{(\columnwidth - 2\tabcolsep) * \real{0.38}}
  >{\raggedright\arraybackslash}p{(\columnwidth - 2\tabcolsep) * \real{0.58}}@{}}
\toprule
Parameter & Interpretation \\
\midrule
\endhead
\(\alpha\) & 初始预期性能 \\
\(\beta\) & 变点前的改善速度 \\
\(\gamma\) & 变点处的斜率变化 \\
\(\beta+\gamma\) & 变点后的改善或退化速度 \\
\(\tau\) & 潜在趋势变化点 \\
\bottomrule
\end{longtable}
\end{readertable}

对于 Better, Then Worse，我们特别关注：
\[\boxed{
\beta>0,
\qquad
\beta+\gamma<0.
}\]

这表示 Agent 在早期具有正向性能趋势，但在后期进入负向性能趋势。

\textbf{Segmented Regression Explorer}

\[S_k^{\mathrm{OOS}}=\alpha+\beta k+\gamma(k-\tau)_+\]

Change point \(\tau\)\\
\textbf{Round 6}

6（范围 2--18）

Initial slope \(\beta\)\\
\textbf{0.045}

0.045（范围 0--0.07）

Slope change \(\gamma\)\\
\textbf{-0.075}

-0.075（范围 -0.12--0.04）

\bookgraphic{0.467}{part-3-48-0001-007-01}{}

Pre-change slope

\textbf{0.045}

Post-change slope

\textbf{-0.030}

\textbf{Better, Then Worse}

确定性教学模型，固定截距 α = 0.4，暂时省略 Evaluation Noise。

tau = 18

\bookgraphic{0.467}{part-3-48-0001-007-state-01}{}

你会发现，仅仅发生斜率变化并不意味着性能退化。

例如：
\[\beta=0.05,\qquad\gamma=-0.02,\]

则：
\[\beta+\gamma=0.03>0.\]

此时 Agent 仍然在改善，只不过改善速度有所下降。

所以 \textbf{Diminishing Returns 不等于 Performance Decay。}

\hypertarget{part-3-48-0001-ux5982ux4f55ux4f30ux8ba1ux672aux77e5ux7684ux53d8ux70b9}{%
\subsubsection{4.3 如何估计未知的变点？}\label{part-3-48-0001-ux5982ux4f55ux4f30ux8ba1ux672aux77e5ux7684ux53d8ux70b9}}

既然我们事先不知道 \(\tau\)，就可以枚举所有合理的候选位置。

对于每一个候选变点 \(\tau\)，通过 Ordinary Least Squares 估计：
\[\hat\alpha_\tau,\quad
\hat\beta_\tau,\quad
\hat\gamma_\tau.\]

并计算 Residual Sum of Squares：
\[\operatorname{RSS}(\tau)
=
\sum_{k=0}^{K}
\left[
S_k^{\mathrm{OOS}}
-\hat\alpha_\tau
-\hat\beta_\tau k
-\hat\gamma_\tau(k-\tau)_+
\right]^2.\]

最后：
\[\boxed{
\hat\tau=
\arg\min_{\tau\in\mathcal T}
\operatorname{RSS}(\tau)
}\]

其中 \(\mathcal T\) 为允许的变点位置集合，需要保证变点两侧均具有足够的观测值。

下面是一个简化的 Python 实现：

\begin{verbatim}
import numpy as np


def estimate_change_point(scores, min_segment=3):
    y = np.asarray(scores, dtype=float)
    k = np.arange(len(y))

    results = []

    for tau in range(
        min_segment,
        len(y) - min_segment
    ):
        X = np.column_stack([
            np.ones(len(y)),
            k,
            np.maximum(k - tau, 0)
        ])

        coefficients, *_ = np.linalg.lstsq(
            X, y, rcond=None
        )

        alpha, beta, gamma = coefficients

        fitted = X @ coefficients
        rss = np.sum((y - fitted) ** 2)

        results.append({
            "change_point": tau,
            "pre_slope": beta,
            "post_slope": beta + gamma,
            "rss": rss
        })

    if not results:
        raise ValueError("Insufficient observations")

    return min(results, key=lambda r: r["rss"])
\end{verbatim}

该代码只是单变点的探索性 OLS Estimator，尚未包含统计检验和置信区间，也未在真实 RSI 数据上运行。

未知变点的分段回归具有成熟的统计研究基础，可以进一步参考 Muggeo (2003) 的工作。

\href{https://doi.org/10.1002/sim.1545}{来源 1}

\hypertarget{part-3-48-0001-ux4e00ux4e2aux975eux5e38ux91cdux8981ux7684ux7edfux8ba1ux9677ux9631}{%
\subsubsection{4.4 一个非常重要的统计陷阱}\label{part-3-48-0001-ux4e00ux4e2aux975eux5e38ux91cdux8981ux7684ux7edfux8ba1ux9677ux9631}}

上述算法只要具有足够的观测值，就一定会返回一个候选变点。

但这并不意味着真实变点一定存在。

因此，需要比较两个竞争模型：
\[\mathcal M_0:
\quad
S_k^{\mathrm{OOS}}
=
\alpha+\beta k+\eta_k,\]

\[\mathcal M_1:
\quad
S_k^{\mathrm{OOS}}
=
\alpha+\beta k+\gamma(k-\tau)_+
+\eta_k.\]

由于 \(\mathcal M_1\) 具有更大的模型灵活性，它的拟合误差通常不会高于 \(\mathcal M_0\)。

所以不能仅仅因为：
\[\operatorname{RSS}_1<\operatorname{RSS}_0\]

就认定存在真正的 Change Point。

还需要考虑模型复杂度、变点搜索产生的选择偏差以及残差相关性。由于在无变点的零假设下 \(\tau\) 不可识别，普通固定参数的 F-test 也不能直接照搬。

对于初期研究，建议采用如下优先级：

\begin{readertable}
\begin{longtable}[]{@{}
  >{\raggedright\arraybackslash}p{(\columnwidth - 4\tabcolsep) * \real{0.40}}
  >{\raggedright\arraybackslash}p{(\columnwidth - 4\tabcolsep) * \real{0.39}}
  >{\raggedright\arraybackslash}p{(\columnwidth - 4\tabcolsep) * \real{0.16}}@{}}
\toprule
Method & Function & Priority \\
\midrule
\endhead
Peak Detection & 描述经验峰值 & P0 \\
Segmented Regression & 估计趋势变化点 & P0 \\
Bootstrap-based Calibration & 估计变点与模型比较的不确定性 & P0 \\
PELT & 检测多个变点 & P1 \\
CUSUM & 检测累计偏离或监控变化 & P1 \\
Bayesian Change-point Detection & 建立变点的概率模型 & P2 \\
\bottomrule
\end{longtable}
\end{readertable}

关于更多 Offline Change-point Detection 方法，可以参考 Truong、Oudre 和 Vayatis (2020) 的综述。

\href{https://www.sciencedirect.com/science/article/pii/S0165168419303494}{来源 1}

现阶段不需要将所有方法都学完。\textbf{Segmented Regression 已经足够作为一个有解释力的研究起点。}

\hypertarget{part-3-48-0001-revision-classificationux4feeux6539ux5206ux7c7b}{%
\subsection{5. Revision Classification（修改分类）}\label{part-3-48-0001-revision-classificationux4feeux6539ux5206ux7c7b}}

识别变点之后，下一步是研究变点附近发生了什么。

\hypertarget{part-3-48-0001-revision-ux4e0dux5e94ux8be5ux53eaux7528ux5355ux4e00ux7c7bux522bux8868ux793a}{%
\subsubsection{5.1 Revision 不应该只用单一类别表示}\label{part-3-48-0001-revision-ux4e0dux5e94ux8be5ux53eaux7528ux5355ux4e00ux7c7bux522bux8868ux793a}}

假设 Agent 执行如下代码修改：

\begin{verbatim}
- candidates = generate_features(n=10)
+ candidates = generate_features(n=100)

- model = Ridge(alpha=1.0)
+ model = Ridge(alpha=0.01)

+ memory.save_best_feature_combinations()
\end{verbatim}

这实际上同时改变了三个组件：

\begin{readertable}
\begin{longtable}[]{@{}
  >{\raggedright\arraybackslash}p{(\columnwidth - 4\tabcolsep) * \real{0.30}}
  >{\raggedright\arraybackslash}p{(\columnwidth - 4\tabcolsep) * \real{0.33}}
  >{\raggedright\arraybackslash}p{(\columnwidth - 4\tabcolsep) * \real{0.32}}@{}}
\toprule
Code Change & Revision Label & Statistical Implication \\
\midrule
\endhead
10 → 100 candidates & Search Expansion & 更强的选择压力 \\
Ridge penalty 1 → 0.01 & Regularization Reduction & 模型约束降低 \\
Save best combinations & Memory Expansion & 历史选择信息积累 \\
\bottomrule
\end{longtable}
\end{readertable}

因此，Revision Classification 更适合被设计为一个 \textbf{Multi-label Classification Problem}，而不是强制每一轮只能归入一个类别。

\hypertarget{part-3-48-0001-proposed-revision-taxonomy}{%
\subsubsection{5.2 Proposed Revision Taxonomy}\label{part-3-48-0001-proposed-revision-taxonomy}}

\begin{readertable}
\begin{longtable}[]{@{}
  >{\raggedright\arraybackslash}p{(\columnwidth - 4\tabcolsep) * \real{0.14}}
  >{\raggedright\arraybackslash}p{(\columnwidth - 4\tabcolsep) * \real{0.37}}
  >{\raggedright\arraybackslash}p{(\columnwidth - 4\tabcolsep) * \real{0.45}}@{}}
\toprule
Code & Category & 修改对象 \\
\midrule
\endhead
R01 & Feature Engineering & 特征构造、筛选、变换 \\
R02 & Model Architecture & 预测模型结构 \\
R03 & Hyperparameter & 模型参数设置 \\
R04 & Search Strategy & 搜索宽度、深度与路径 \\
R05 & Memory \& Context & 历史信息保存与检索 \\
R06 & Evaluation Logic & 评价指标及验证协议 \\
R07 & Data Processing & 数据清洗、划分与时间处理 \\
R08 & Regularization & 复杂度及稳健性约束 \\
R09 & Resource Allocation & 时间、Token、候选预算 \\
R10 & Infrastructure & 调试、并行化、执行优化 \\
R11 & Improver Logic & 自我改进机制 \\
\bottomrule
\end{longtable}
\end{readertable}

值得注意的是，Revision Category 与 Failure Mechanism 并不是同一个概念。

例如，Feature Expansion 可能产生真实改进，也可能引发过拟合。分类应当先描述发生了什么修改，再独立调查其统计影响。

\hypertarget{part-3-48-0001-revision-impact-quadrants}{%
\subsubsection{5.3 Revision Impact Quadrants}\label{part-3-48-0001-revision-impact-quadrants}}

对于每一次被接受的 Revision，可以定义：
\[\Delta_k^{\mathrm{IS}},
\qquad
\Delta_k^{\mathrm{OOS}}.\]

将每次修改放入一个二维坐标系中。

\textbf{Revision Impact Map}

\bookgraphic{0.350}{part-3-48-0001-010-01}{假想 Revision 数据。右下象限是重点诊断区域，但进入该区域并不能单独证明过拟合。}

我们可以据此定义一个描述性统计量：
\[\widehat p_r
=
\frac{
\sum_k
\mathbf 1_{\{
r\in\mathcal R_k,\,
\Delta_k^{\mathrm{IS}}>0,\,
\Delta_k^{\mathrm{OOS}}<0
\}}
}{
\sum_k\mathbf 1_{\{r\in\mathcal R_k\}}
},\]

其中：

\begin{itemize}
\tightlist
\item
  \(r\)：某种 Revision Category；
\item
  \(\mathcal R_k\)：第 \(k\) 次修改包含的类别集合。
\end{itemize}

它衡量包含某类修改的实验中，有多少比例表现为 IS 改善、OOS 下降。

但这个比例仅仅是 Association，并非该类修改导致过拟合的因果概率。不同修改可能共同出现，而且后期修改本身就可能处在更容易发生退化的阶段。

\hypertarget{part-3-48-0001-rollback-counterfactual-and-ablation-study}{%
\subsection{6. Rollback, Counterfactual and Ablation Study}\label{part-3-48-0001-rollback-counterfactual-and-ablation-study}}

这是从描述性分析进入机制验证的关键一步。

\hypertarget{part-3-48-0001-rollbackux56deux6edaux5b9eux9a8c}{%
\subsubsection{6.1 Rollback（回滚实验）}\label{part-3-48-0001-rollbackux56deux6edaux5b9eux9a8c}}

假设某个 Agent 的演化过程如下：
\[A_0\rightarrow A_1\rightarrow
A_2\rightarrow A_3\rightarrow A_4.\]

我们发现 \(A_3\) 的 OOS Performance 明显下降。

一种自然的方法是回滚到 \(A_2\)，重新执行实验。

但回滚实际上可以回答两种不同的问题。

\begin{readertable}
\begin{longtable}[]{@{}
  >{\raggedright\arraybackslash}p{(\columnwidth - 2\tabcolsep) * \real{0.34}}
  >{\raggedright\arraybackslash}p{(\columnwidth - 2\tabcolsep) * \real{0.62}}@{}}
\toprule
Experiment & Question \\
\midrule
\endhead
Static Rollback & 旧版本是否具有更好的当前性能？ \\
Evolutionary Rollback & 从旧版本继续改进，是否能够形成更好的长期轨迹？ \\
\bottomrule
\end{longtable}
\end{readertable}

前者关注版本的直接性能，后者关注 Path Dependence 和未来演化潜力。

两者不能混为一谈。

\hypertarget{part-3-48-0001-counterfactual-revision-analysisux53cdux4e8bux5b9eux4feeux6539ux5206ux6790}{%
\subsubsection{6.2 Counterfactual Revision Analysis（反事实修改分析）}\label{part-3-48-0001-counterfactual-revision-analysisux53cdux4e8bux5b9eux4feeux6539ux5206ux6790}}

假设 \(A_3\) 相对于 \(A_2\) 包含三个修改：
\[A_3=A_2+\{r_1,r_2,r_3\}.\]

这里的加号表示应用一组代码修改，并非通常意义上的代数运算。

为了分析 \(r_2\) 的影响，我们构造：
\[A_3^{(-r_2)}
=
A_2+\{r_1,r_3\}.\]

然后比较：
\[\boxed{
\delta_{r_2}^{\mathrm{OOS}}
=
J^{\mathrm{OOS}}(A_3)
-
J^{\mathrm{OOS}}(A_3^{(-r_2)})
}\]

其中 \(\delta_{r_2}^{\mathrm{OOS}}\) 表示在其余修改保持不变时，保留 \(r_2\) 所对应的性能差异。

\bookgraphic{0.439}{part-3-48-0001-011-01}{假想实验结果。只有在其余组件、资源预算及评价条件得到充分控制时，这种差异才具有明确的干预解释。}

这比观察 Revision 前后的得分变化更具有解释力，因为我们主动改变了一个实验因素。

\hypertarget{part-3-48-0001-single-component-ablation-ux7684ux5c40ux9650}{%
\subsubsection{6.3 Single-component Ablation 的局限}\label{part-3-48-0001-single-component-ablation-ux7684ux5c40ux9650}}

但删除单个组件也存在一个问题：不同修改可能具有 Interaction Effect。

例如：

\begin{readertable}
\begin{longtable}[]{@{}
  >{\raggedright\arraybackslash}p{(\columnwidth - 2\tabcolsep) * \real{0.70}}
  >{\raggedleft\arraybackslash}p{(\columnwidth - 2\tabcolsep) * \real{0.26}}@{}}
\toprule
Configuration & OOS Score \\
\midrule
\endhead
Baseline & 0.60 \\
Feature Expansion & 0.50 \\
Regularization & 0.63 \\
Feature Expansion + Regularization & 0.72 \\
\bottomrule
\end{longtable}
\end{readertable}

这里 Feature Expansion 单独使用时损害性能，但与 Regularization 结合之后反而能够获得更好的结果。

因此不能简单认定 Feature Expansion 是一个有害的 Revision Category。

对于高度相关的组件，可以进一步采用 III-B 中介绍过的 Factorial Experiment。

\hypertarget{part-3-48-0001-direct-effect-vs.-evolutionary-effect}{%
\subsubsection{6.4 Direct Effect vs. Evolutionary Effect}\label{part-3-48-0001-direct-effect-vs.-evolutionary-effect}}

这一点是 RSI 特有的研究难题。

某次 Revision 可能只使当前 OOS Score 提高 0.01，但却显著改变后续候选生成机制，使未来十轮的表现大幅提高或下降。

因此我们可以分别研究：
\[\delta_r^{\mathrm{direct}}
=
J^{\mathrm{OOS}}(A^+)
-
J^{\mathrm{OOS}}(A^-),\]

以及：
\[\boxed{
\delta_r^{\mathrm{evolution}}(L)
=
\mathbb E\left[
J^{\mathrm{OOS}}(A_L^+)
-
J^{\mathrm{OOS}}(A_L^-)
\right].
}\]

其中 \(A^+\) 和 \(A^-\) 分别表示包含与不包含该修改的版本；\(A_L^+\) 和 \(A_L^-\) 表示从两个版本出发，在匹配预算下继续执行 \(L\) 轮改进得到的 Agent。

后者衡量的不是某次修改的即时收益，而是其对未来演化过程的影响。

\textbf{A Revision can improve the agent while damaging its future evolution.}

这可能是 RSI Diagnosis 区别于普通 Model Ablation 的重要研究方向。

\hypertarget{part-3-48-0001-statistical-analysis-across-independent-runs}{%
\subsection{7. Statistical Analysis across Independent Runs}\label{part-3-48-0001-statistical-analysis-across-independent-runs}}

\hypertarget{part-3-48-0001-ux4e3aux4ec0ux4e48ux4e00ux6761ux6f02ux4eaeux7684ux66f2ux7ebfux8fdcux8fdcux4e0dux591f}{%
\subsubsection{7.1 为什么一条漂亮的曲线远远不够？}\label{part-3-48-0001-ux4e3aux4ec0ux4e48ux4e00ux6761ux6f02ux4eaeux7684ux66f2ux7ebfux8fdcux8fdcux4e0dux591f}}

LLM-based RSI 通常具有较强随机性。

相同的初始 Agent、相同的任务和相同的 Evaluation Protocol，也可能产生完全不同的演化路径。

\bookgraphic{0.723}{part-3-48-0001-012-01}{}

假想的四次独立 RSI Runs：不同实验具有不同的退化时间，甚至可能完全不发生退化。

如果我们只展示 Run 1 和 Run 3，就可能夸大 RSI 退化现象的普遍性。

如果只展示 Run 4，又可能掩盖 RSI 的统计风险。

因此，需要报告所有符合预先规定实验条件的 Runs，包括失败、中断和未检测到退化的实验。

\hypertarget{part-3-48-0001-mean-trajectory-ux4e0e-uncertainty}{%
\subsubsection{7.2 Mean Trajectory 与 Uncertainty}\label{part-3-48-0001-mean-trajectory-ux4e0e-uncertainty}}

设我们执行了 \(N\) 次独立实验，其中 \(S_{r,k}^{\mathrm{OOS}}\) 表示第 \(r\) 次实验、第 \(k\) 轮的 OOS 得分。

定义平均轨迹：
\[\boxed{
\bar S_k^{\mathrm{OOS}}
=
\frac1N\sum_{r=1}^{N}
S_{r,k}^{\mathrm{OOS}}
}\]

在 Runs 独立同分布且具有有限方差的条件下，对平均得分的标准误差可以估计为：
\[\widehat{\operatorname{SE}}_k
=
\frac{s_k}{\sqrt N},\]

其中 \(s_k\) 为第 \(k\) 轮不同 Runs 之间的样本标准差。

\textbf{Mean Trajectory with Uncertainty}

\bookgraphic{0.465}{part-3-48-0001-014-01}{示意性平均轨迹及不确定性区间，并非根据上图四组数据计算的实际置信区间。}

但是，同一次实验中相邻版本的性能显然具有较强相关性。

因此，不能把每个 Iteration 都当成独立观测值。

比较完整演化过程时，应当尽可能以 \textbf{Run} 为重采样或统计比较单位，并保留每条轨迹内部的依赖关系。

还需要区分：

\begin{readertable}
\begin{longtable}[]{@{}
  >{\raggedright\arraybackslash}p{(\columnwidth - 2\tabcolsep) * \real{0.43}}
  >{\raggedright\arraybackslash}p{(\columnwidth - 2\tabcolsep) * \real{0.53}}@{}}
\toprule
Uncertainty Source & Meaning \\
\midrule
\endhead
Agent Stochasticity & LLM 采样和搜索随机性 \\
Evaluation Uncertainty & 有限 OOS 样本造成的得分波动 \\
Environment Variation & 不同市场时期和任务分布的差异 \\
\bottomrule
\end{longtable}
\end{readertable}

如果所有 Runs 都使用同一套金融数据，那么跨 Run 的置信区间主要描述算法随机性，并不充分覆盖未来市场的不确定性。

\hypertarget{part-3-48-0001-decay-onset-ux672cux8eabux4e5fux662fux968fux673aux53d8ux91cf}{%
\subsubsection{7.3 Decay Onset 本身也是随机变量}\label{part-3-48-0001-decay-onset-ux672cux8eabux4e5fux662fux968fux673aux53d8ux91cf}}

设第 \(r\) 次独立实验检测到的变点为：
\[\hat\tau_r.\]

我们不仅可以研究平均性能轨迹，还可以研究：
\[\hat\tau_1,\ldots,\hat\tau_N\]

的分布。

这会产生几个新的 Research Metrics：

\begin{readertable}
\begin{longtable}[]{@{}
  >{\raggedright\arraybackslash}p{(\columnwidth - 2\tabcolsep) * \real{0.38}}
  >{\raggedright\arraybackslash}p{(\columnwidth - 2\tabcolsep) * \real{0.58}}@{}}
\toprule
Metric & Research Meaning \\
\midrule
\endhead
Decay Incidence & 有多少比例的实验检测到了退化 \\
Median Decay Onset & 退化通常何时开始 \\
Decay Severity & 退化发生后有多严重 \\
Time-to-decay & 在多大评估预算之后开始退化 \\
\bottomrule
\end{longtable}
\end{readertable}

对于直到实验结束仍未发生退化的 Runs，可以考虑将其视为 Right-censored Observations，引入 Survival Analysis。

不过，Survival Analysis 暂时属于进阶内容，而且必须先建立可靠的退化事件定义。检测结果本身不可靠时，更复杂的分析模型也不能自动解决问题。

\hypertarget{part-3-48-0001-reproducibilityux53efux590dux73b0ux6027}{%
\subsection{8. Reproducibility（可复现性）}\label{part-3-48-0001-reproducibilityux53efux590dux73b0ux6027}}

RSI 的复现比普通机器学习实验更加复杂，因为它的最终结果不仅取决于模型参数，还取决于整个演化历史。

\[A_{k+1}
\sim
\mathcal K(\cdot\mid A_k,H_k).\]

因此，保存 \(A_K\) 并不足以完整复现整个过程。

\hypertarget{part-3-48-0001-three-levels-of-reproducibility}{%
\subsubsection{8.1 Three Levels of Reproducibility}\label{part-3-48-0001-three-levels-of-reproducibility}}

\begin{readertable}
\begin{longtable}[]{@{}
  >{\raggedright\arraybackslash}p{(\columnwidth - 2\tabcolsep) * \real{0.43}}
  >{\raggedright\arraybackslash}p{(\columnwidth - 2\tabcolsep) * \real{0.53}}@{}}
\toprule
Level & Research Goal \\
\midrule
\endhead
Artifact Reproduction & 重新运行某个已经生成的 Agent \\
Trajectory Replay & 重建此前发生的修改和选择过程 \\
Statistical Replication & 独立运行系统并复现总体统计结论 \\
\bottomrule
\end{longtable}
\end{readertable}

第三层通常比要求每一次 LLM 输出完全相同更具有科研价值。

尤其是在使用商业 LLM API 时，即使指定了相同的 Seed，也未必能够保证输出完全确定。

\hypertarget{part-3-48-0001-reproducibility-checklist}{%
\subsubsection{8.2 Reproducibility Checklist}\label{part-3-48-0001-reproducibility-checklist}}

\textbf{Research Reproducibility Checklist}

\begin{itemize}
\item
  固定并记录 Foundation Model 及 API 配置
\item
  保存初始 Agent 和全部候选代码版本
\item
  记录 Prompt、Memory 和工具调用历史
\item
  保存数据版本、划分方法与文件校验值
\item
  固定外部 Evaluator 和评分协议
\item
  记录随机种子及依赖环境
\item
  记录全部成功、失败和被拒绝的实验
\item
  记录 Token、Runtime 和 Evaluation Budget
\item
  隔离 OOS Audit 与 Agent Development
\item
  保留能够复现实验结果的分析脚本
\end{itemize}

适用于未来实际搭建或审阅 RSI Experimental Framework。

你的实验日志不应该仅仅服务于论文写作，而应该能够成为整个研究结论的 Evidence Chain。

\hypertarget{part-3-48-0001-mitigation-mechanismsux7f13ux89e3ux673aux5236}{%
\subsection{9. Mitigation Mechanisms（缓解机制）}\label{part-3-48-0001-mitigation-mechanismsux7f13ux89e3ux673aux5236}}

现在我们可以将 Part II 和 Part III 中已经讨论过的干预手段统一整理。

\begin{readertable}
\begin{longtable}[]{@{}
  >{\raggedright\arraybackslash}p{(\columnwidth - 4\tabcolsep) * \real{0.38}}
  >{\raggedright\arraybackslash}p{(\columnwidth - 4\tabcolsep) * \real{0.21}}
  >{\raggedright\arraybackslash}p{(\columnwidth - 4\tabcolsep) * \real{0.38}}@{}}
\toprule
Mitigation & 干预位置 & Targeted Failure \\
\midrule
\endhead
Early Stopping & Outer Loop & 后期过拟合 \\
Search Budget Control & Generation & Multiple Testing \\
Feedback Compression & Evaluator & Information Accumulation \\
Reusable Holdout & Evaluation & Adaptive Overfitting \\
Complexity Penalty & Selection & Complexity Accumulation \\
Statistical Acceptance & Selection & Noise-driven Revision \\
Sparse Modification & Generation & 复杂度累积、归因困难 \\
Memory Filtering & Memory & Spurious Experience \\
Archive-based Search & Controller & Premature Convergence \\
Independent Verification & Evaluation & Leakage、Reward Hacking \\
\bottomrule
\end{longtable}
\end{readertable}

它们并不是可以无条件解决问题的十种方法，而是针对不同 Failure Mechanism 的候选干预。

例如，如果性能退化主要来自 Concept Drift，那么单纯限制 Agent 的 Search Budget 可能没有太大效果。

这说明研究流程应当是：
\[\boxed{
\text{Diagnosis}
\longrightarrow
\text{Targeted Intervention}
\longrightarrow
\text{Independent Verification}
}\]

而不是在还没有理解失败原因时，直接将所有 Regularization Mechanisms 添加到系统中。

\hypertarget{part-3-48-0001-research-practice-ux4e00ux4e2aux5b8cux6574ux7684ux6700ux5c0fux7814ux7a76ux9879ux76ee}{%
\subsection{10. Research Practice --- 一个完整的最小研究项目}\label{part-3-48-0001-research-practice-ux4e00ux4e2aux5b8cux6574ux7684ux6700ux5c0fux7814ux7a76ux9879ux76ee}}

假设现在我们希望独立开展一个与 X Program \#161 相关的 Pilot Study。

我建议先构建一个不涉及昂贵 LLM 调用的可控实验，再逐步引入真实 Agent。

\textbf{Phase 1 --- Synthetic Financial Environment}

\begin{verbatim}
  构造包含真实预测信号与无效噪声特征的数据集，控制 SNR。
\end{verbatim}

\textbf{Phase 2 --- Non-adaptive Baseline}

\begin{verbatim}
  通过 Random Search 建立普通 Model Selection 的基线。
\end{verbatim}

\textbf{Phase 3 --- Adaptive Optimization}

\begin{verbatim}
  允许系统根据历史 IS Feedback 调整特征搜索方向。
\end{verbatim}

\textbf{Phase 4 --- Self-modifying Optimizer}

\begin{verbatim}
  允许优化机制本身改变，包括候选生成、搜索宽度与选择规则。
\end{verbatim}

\textbf{Phase 5 --- Diagnosis}

\begin{verbatim}
  记录完整轨迹，执行 Segmented Regression、Revision Classification 和 Rollback。
\end{verbatim}

\textbf{Phase 6 --- Mitigation}

\begin{verbatim}
  引入复杂度约束、反馈控制或统计接受机制，在独立环境中检验效果。
\end{verbatim}

这项研究最有价值的地方是，我们可以逐步增加系统的复杂度。

如果仅仅通过 Phase 2 就能够复现 Better, Then Worse，那么不应该将全部退化归因于 RSI。

如果 Phase 3 明显加重退化，则说明 Adaptive Feedback 值得关注。

如果在严格控制预算之后，Phase 4 仍然表现出额外的退化，就有理由继续研究 Recursive Self-Modification 的独特影响。

这样的实验结果即使不支持最初的研究假设，也仍然具有科学价值。

\hypertarget{part-3-48-0001-literature-index}{%
\subsection{11. Literature Index}\label{part-3-48-0001-literature-index}}

本阶段不需要额外阅读大量 Agent 论文，主要补充两篇统计方法文献。

\textbf{P0 --- Estimating Regression Models with Unknown Break-points}

Vito M. R. Muggeo · Statistics in Medicine · 2003 · 同行评审

重点学习未知变点的分段回归建模及估计方法。

\href{https://doi.org/10.1002/sim.1545}{DOI: 10.1002/sim.1545} \href{https://pubmed.ncbi.nlm.nih.gov/12973787/}{来源 1}

\textbf{P1 --- Selective Review of Offline Change Point Detection Methods}

Charles Truong, Laurent Oudre, Nicolas Vayatis · Signal Processing · 2020 · 同行评审

系统了解 Change-point Detection 中的 Cost Function、Search Method 和模型选择问题。

\href{https://www.sciencedirect.com/science/article/pii/S0165168419303494}{正式出版页面} \href{https://www.sciencedirect.com/science/article/pii/S0165168419303494}{来源 1}

上述两篇论文是统计方法参考，并不是专门针对 RSI 的研究。如何将这些方法用于短而高度相关的 Agent Trajectories，仍需要根据实际实验条件进行适配。

\hypertarget{part-3-48-0001-iii-f-conclusion-index}{%
\subsection{12. III-F --- Conclusion Index}\label{part-3-48-0001-iii-f-conclusion-index}}

\begin{readertable}
\begin{longtable}[]{@{}
  >{\raggedright\arraybackslash}p{(\columnwidth - 4\tabcolsep) * \real{0.12}}
  >{\raggedright\arraybackslash}p{(\columnwidth - 4\tabcolsep) * \real{0.69}}
  >{\raggedright\arraybackslash}p{(\columnwidth - 4\tabcolsep) * \real{0.15}}@{}}
\toprule
ID & Key Conclusion & Index \\
\midrule
\endhead
F1 & OOS Decay 不等于已经确认的 Overfitting & §1 \\
F2 & Instrumentation 必须保存完整搜索历史 & §2 \\
F3 & Proposal Count、Lineage Depth 和 Evaluation Budget 不等价 & §3.1 \\
F4 & Endpoint Improvement 可能掩盖严重的中间退化 & §3.2 \\
F5 & Observed Peak 不等于真实 Change Point & §4.1 \\
F6 & Segmented Regression 可以估计趋势变化 & §4.2--3 \\
F7 & 拟合出一个变点不代表真实变点存在 & §4.4 \\
F8 & Revision Classification 应当允许多标签 & §5 \\
F9 & Revision Association 不等于因果效应 & §5.3 \\
F10 & Rollback 和 Ablation 有助于验证具体修改的影响 & §6 \\
F11 & Direct Effect 与 Evolutionary Effect 需要区分 & §6.4 \\
F12 & 独立重复实验是统计可靠性的基础 & §7 \\
F13 & RSI Reproducibility 需要保存完整演化历史 & §8 \\
F14 & Mitigation 应当针对已经识别的失败机制 & §9 \\
\bottomrule
\end{longtable}
\end{readertable}

\hypertarget{part-3-48-0001-part-iii-consolidated-conclusion-index}{%
\subsection{Part III --- Consolidated Conclusion Index}\label{part-3-48-0001-part-iii-consolidated-conclusion-index}}

至此，我们正式完成了 Part III。

按照之前的学习方式，这里不再重复全部推导，而是将六个模块的本质结论统一整理。

\hypertarget{part-3-48-0001-core-knowledge-map}{%
\subsection{13. Core Knowledge Map}\label{part-3-48-0001-core-knowledge-map}}

\begin{readertable}
\begin{longtable}[]{@{}
  >{\raggedright\arraybackslash}p{(\columnwidth - 4\tabcolsep) * \real{0.12}}
  >{\raggedright\arraybackslash}p{(\columnwidth - 4\tabcolsep) * \real{0.38}}
  >{\raggedright\arraybackslash}p{(\columnwidth - 4\tabcolsep) * \real{0.46}}@{}}
\toprule
Module & Fundamental Question & Core Knowledge \\
\midrule
\endhead
III-A & Agent 是什么？ & Architecture、Harness、State \\
III-B & Agent 如何执行任务和积累经验？ & Tools、Memory、Planning、Reflection \\
III-C & Agent 如何递归改进自身？ & Self-reference、Meta-optimization \\
III-D & Agent 如何选择演化方向？ & Search、Evaluation、Selection Pressure \\
III-E & 真实 RSI 系统如何实现？ & STOP、AIDE、ADAS、DGM 等 \\
III-F & 如何研究 RSI 的性能退化？ & Instrumentation、Change Point、Ablation \\
\bottomrule
\end{longtable}
\end{readertable}

\hypertarget{part-3-48-0001-core-mathematical-results}{%
\subsection{14. Core Mathematical Results}\label{part-3-48-0001-core-mathematical-results}}

\hypertarget{part-3-48-0001-result-1-agent-representation}{%
\paragraph{Result 1 --- Agent Representation}\label{part-3-48-0001-result-1-agent-representation}}

\[A=(\mathcal L,P,M,T,C).\]

Agent 的能力由底层模型与外围软件系统共同决定。

Index：III-A §3

\hypertarget{part-3-48-0001-result-2-agent-execution}{%
\paragraph{Result 2 --- Agent Execution}\label{part-3-48-0001-result-2-agent-execution}}

\[\begin{cases}
a_t\sim\pi_A(\cdot\mid s_t),\\
(e_{t+1},o_{t+1})
\sim\mathcal E(\cdot\mid e_t,a_t),\\
s_{t+1}=U_A(s_t,a_t,o_{t+1}).
\end{cases}\]

Agent 可以被抽象为与环境持续交互的状态转换系统。

Index：III-A §3.2

\hypertarget{part-3-48-0001-result-3-memory-mediated-adaptation}{%
\paragraph{Result 3 --- Memory-mediated Adaptation}\label{part-3-48-0001-result-3-memory-mediated-adaptation}}

\[M_{k+1}=U_M(M_k,F_k,H_k).\]

Memory 可以积累可迁移经验，也可能积累开发数据中特有的统计偏差。

Index：III-B §3

\hypertarget{part-3-48-0001-result-4-recursive-meta-optimization}{%
\paragraph{Result 4 --- Recursive Meta-optimization}\label{part-3-48-0001-result-4-recursive-meta-optimization}}

\[\boxed{
\mathcal I_{k+1}
=
\mathcal I_k
(\widehat J_D,\mathcal I_k,\mathcal L)
}\]

当前优化器同时充当优化主体与待优化对象。

Index：III-C §3

\hypertarget{part-3-48-0001-result-5-path-dependent-evolution}{%
\paragraph{Result 5 --- Path-dependent Evolution}\label{part-3-48-0001-result-5-path-dependent-evolution}}

\[A_{k+1}\sim\mathcal K(\cdot\mid A_k,H_k).\]

未来演化受到当前 Agent 及历史信息的共同影响。

Index：III-C §6

\hypertarget{part-3-48-0001-result-6-candidate-decay-model}{%
\paragraph{Result 6 --- Candidate Decay Model}\label{part-3-48-0001-result-6-candidate-decay-model}}

\[\begin{cases}
S_k^{\mathrm{IS}}=S_0+B_k+V_k,\\
S_k^{\mathrm{OOS}}=S_0+B_k-\lambda V_k.
\end{cases}\]

在这一解释性模型中，当真实边际收益不足以抵消过度适应造成的损害时，OOS Performance 开始下降。

\[\Delta B_k<\lambda\Delta V_k.\]

Index：III-C §7.3。属于假设性 Toy Model，而非已验证的 RSI 定律。

\hypertarget{part-3-48-0001-result-7-greedy-selection}{%
\paragraph{Result 7 --- Greedy Selection}\label{part-3-48-0001-result-7-greedy-selection}}

\[A_{k+1}
=
\arg\max_{A\in\mathcal C_k\cup\{A_k\}}
\widehat J_{\mathrm{IS}}(A).\]

固定、确定性的 IS Score 可以在这一规则下机械地保持单调增长。

Index：III-D §2.1

\hypertarget{part-3-48-0001-result-8-regularized-agent-optimization}{%
\paragraph{Result 8 --- Regularized Agent Optimization}\label{part-3-48-0001-result-8-regularized-agent-optimization}}

\[U(A)
=
\widehat J_{\mathrm{IS}}(A)
-\lambda_1\Omega(A)
-\lambda_2C_{\mathrm{res}}(A).\]

通过复杂度与资源约束限制无节制的 Proxy Optimization。

Index：III-D §6

\hypertarget{part-3-48-0001-result-9-performance-change-point-model}{%
\paragraph{Result 9 --- Performance Change-point Model}\label{part-3-48-0001-result-9-performance-change-point-model}}

\[\boxed{
S_k^{\mathrm{OOS}}
=
\alpha+\beta k+\gamma(k-\tau)_++\eta_k
}\]

在所设模型下，Better, Then Worse 对应：
\[\beta>0,\qquad\beta+\gamma<0.\]

Index：III-F §4

\hypertarget{part-3-48-0001-result-10-evolutionary-revision-effect}{%
\paragraph{Result 10 --- Evolutionary Revision Effect}\label{part-3-48-0001-result-10-evolutionary-revision-effect}}

\[\boxed{
\delta_r^{\mathrm{evolution}}(L)
=
\mathbb E[
J^{\mathrm{OOS}}(A_L^+)
-
J^{\mathrm{OOS}}(A_L^-)
]
}\]

研究特定修改对后续演化过程的长期影响，而不仅是即时性能变化。

Index：III-F §6.4

\hypertarget{part-3-48-0001-part-ii-part-iii-unified-research-framework}{%
\subsection{15. Part II + Part III --- Unified Research Framework}\label{part-3-48-0001-part-ii-part-iii-unified-research-framework}}

现在我们可以将之前两大部分真正整合起来。

\textbf{Part II --- Statistics}

Finite-sample Error

Selection Bias

Adaptive Overfitting

Distribution Shift

\textbf{Part III --- Systems}

Agent Harness

Recursive Optimization

Search \& Selection

Memory \& Feedback

↓

\textbf{Statistical Bias Propagating through Recursive Agent Evolution}

统计偏差通过 Agent 的搜索、记忆与递归修改机制不断传播

↓

\textbf{Performance Trajectory}

IS Growth · OOS Improvement · Potential Decay

↓

\textbf{Detection}

When?

\textbf{Diagnosis}

Why?

↓

\textbf{Mitigation \& Independent Verification}

How can we improve generalization?

学到这里，我们已经能够从统计学和计算机科学两个角度解释该课题。

不过，目前仍然缺少第三个关键维度：

\textbf{为什么 Finance 特别容易成为 RSI Overfitting 的高风险环境？}

前面所有的统计推导，很多都是建立在 IID、有限方差、有界损失或相对稳定的数据生成机制之上。

真实金融数据恰恰经常不满足这些条件。

这将是 Part IV 的核心。

\hypertarget{part-3-48-0001-next-part-iv-financial-ai-and-quantitative-research}{%
\subsection{Next --- Part IV: Financial AI and Quantitative Research}\label{part-3-48-0001-next-part-iv-financial-ai-and-quantitative-research}}

Part I --- Understanding the Research Problem

Completed

Part II --- Statistical Learning Theory

Completed

Part III --- Agentic AI \& Recursive Self-Improvement

Completed

Part IV --- Financial AI \& Quantitative Research

Next

Part V --- Literature \& Research Frontier

Part VI --- Research Practice

\hypertarget{part-3-48-0001-part-iv-learning-roadmap}{%
\subsection{Part IV --- Learning Roadmap}\label{part-3-48-0001-part-iv-learning-roadmap}}

\begin{readertable}
\begin{longtable}[]{@{}
  >{\raggedright\arraybackslash}p{(\columnwidth - 4\tabcolsep) * \real{0.17}}
  >{\raggedright\arraybackslash}p{(\columnwidth - 4\tabcolsep) * \real{0.61}}
  >{\raggedright\arraybackslash}p{(\columnwidth - 4\tabcolsep) * \real{0.18}}@{}}
\toprule
Section & Topic & Priority \\
\midrule
\endhead
IV-A & Financial Time Series \& Data Generating Process & P0 \\
IV-B & Low Signal-to-Noise Ratio \& Predictability & P0 \\
IV-C & Financial Factors \& Machine Learning Models & P1 \\
IV-D & Returns, Sharpe Ratio \& Performance Evaluation & P0 \\
IV-E & Backtest Overfitting, PBO \& Deflated Sharpe Ratio & P0 \\
IV-F & Market Regime, Non-stationarity \& Robustness & P0 \\
\bottomrule
\end{longtable}
\end{readertable}

我们暂时不深入 Options Pricing、Stochastic Calculus 和复杂的金融衍生品定价，因为这些并不是当前研究所必需的知识。

此外，目前仍然没有获得课题组的实际 Framework、Dataset 和具体 Financial Task。因此，本阶段提出的诊断指标、日志结构、变点算法和实验方案均属于待验证的研究设计。

\textbf{Part II 解决了 Why Overfitting Happens；Part III 解决了 How RSI Amplifies or Mitigates It；Part IV 将解决 Why Finance Makes This Problem Particularly Difficult。}
